\documentclass[../main.tex]{subfiles}

\begin{document}

\section{LITERATURE REVIEW}


 \begin{itemize}
        \item {  \textbf{Reviewing systematic studies on IoT and AI in poultry farming:}}
        There is a systematic review of the literature by Ojo et al.That provides a comprehensive description of AI-based IoT systems in poultry welfare and health management. The review poses challenges in the current poultry management, such as the lack of uniform parameters for welfare assessment and robust monitoring systems to prevent diseases, particularly in broilers. It describes the potential of digital technology in facilitating automation of poultry activities, stimulating high-quality and low-cost production. The study identifies significant uses for example in real-time monitoring of the environment, behavior, and disease detection necessary to improve poultry welfare and productivity.

        \item {\textbf{Discussing IoT applications in behavioral monitoring:}}
     A specific instance of IoT application in poultry rearing is tracking chicken behavior to assess health condition. Ahmed et al.It proposed an IoT-based smart monitoring system applying sensors to detect and analyze chicken behavior, such as feeding, preening, and dustbathing. Applying machine learning algorithms, Random Forest, the system achieves a 98\% accuracy in predicting poultry behaviors, which can indicate health issues. The system uses a better Synthetic Minority Over-sampling Technique (SMOTE) in an Artificial Hummingbird Algorithm (AHA) to manage data imbalance with 97\% accuracy. The approach allows for early warnings on abnormalities to enable timely intervention to prevent disease outbreaks and enhance farm management.

    
     \item {\textbf{ Discussing IoT applications in behavioral monitoring:}}
     IoT solutions have been developed in the management of hatcheries to automate incubation processes. For example, PoultryMon, an IoT-based device, provides live temperature and humidity monitoring of incubators to farmers with critical information to enable them to maintain conditions in top shape . The system connects various incubators, collects data on hatchery and farm scales, and provides outlier alerts, with the ability to increase efficiency by 3\%. These technologies are crucial in ensuring high hatch rates and healthy chicks, addressing challenges like manual monitoring that can lead to production inefficiencies.
    
      \item {\textbf{Examining IoT for sustainability in poultry farming:}}
    Besides, Kamor et al. It examined the use of IoT to make poultry farming more sustainable, with emphasis placed on the Indian context. The paper reports on the development and design of an IoT-based system that incorporates sensors, wireless communication, and cloud-based data processing to monitor and control environmental conditions. The system addresses issues such as process-intensive methodologies, environmental issues, and disease management, resulting in more effective production efficiency and sustainability. The authors state that world chicken meat production has grown to 102 million metric tons in 2020, highlighting the scale and importance of such innovations.

\item {\textbf{Highlighting environmental monitoring systems:}}
Lashari et al. proposed an IoT-based framework for sustainable environmental management of big indoor complexes, including poultry farms. Their framework monitors temperature, humidity, oxygen, and gas levels (e.g., ammonia, CO$_2$, CO) to maintain perfect conditions for healthy birds. Their framework is cost-effective and robust, which may be adopted for large-scale units suffering from environmental problems such as global warming.


\item {\textbf{Discussing cost-effective IoT solutions:}}
Karun et al. created a smart poultry management system with embedded frameworks and mobile technology, emphasizing cost-saving and quality-oriented methods. The system is automated for water, feed, and environmental management through sensors and cloud data accessed on a mobile platform. This solution best suits the needs for addressing labor concerns and optimal productivity in small to medium-scale farms.


\item {\textbf{ Identifying research gaps and opportunities:}}
While these advances have been achieved, the literature identifies a few areas that require further research. One such area is the technical issue of handling large amounts of data generated by IoT devices and converging multiple sensors. Furthermore, ensuring that these technologies are within reach and cost-effective for small farmers, particularly in developing countries such as Nepal, is critical to their mass usage. The lack of localized research focused on realizing the environmental and economic context of Nepal suggests a gap that the IOT project intends to fill.

\item {\textbf{Summarizing the role of IoT in poultry farming:}}
In short, the literature reveals that IoT technology has a huge potential to transform poultry farming with the provision of real-time monitoring, automation, and data-driven decision-making. The IOT project builds upon these grounds in an effort to develop an end-to-end IoT-based system tailored to poultry hatchery and farm management within the Nepalese scenario, addressing global concerns as well as local needs.


    \end{itemize}
\end{document}